% AI usage log for Section 3 of the HW4 write-up.
% Import directly into your report with:
%   % AI usage log for Section 3 of the HW4 write-up.
% Import directly into your report with:
%   % AI usage log for Section 3 of the HW4 write-up.
% Import directly into your report with:
%   % AI usage log for Section 3 of the HW4 write-up.
% Import directly into your report with:
%   \input{ai_usage_log.tex}
% Requires \usepackage{enumitem} if you want to keep the [leftmargin=*] styling below;
% otherwise just remove that option.

\begin{itemize}[leftmargin=*]

\item \textbf{Directory traversal / dataset indexing.}
I asked how to enumerate the per-class image folders under \texttt{hw4-data/} and build
an index of (path, label) pairs for \texttt{SUN397Dataset}. The AI pointed me to
\texttt{os.listdir}/\texttt{os.walk}/\texttt{os.path.join}/\texttt{sorted()} conceptually
but did not write the loop. I initially used \texttt{os.walk(dir\_path)} without consuming
it as a generator (tried to unpack it directly into three names), which the AI caught by
asking me to reconsider what a generator yields. I also hit a \texttt{NotADirectoryError}
from \texttt{.DS\_Store} and from checking \texttt{os.path.isfile(dir)} on the wrong
(un-joined) path; both were my bugs, caught via tracebacks I pasted and the AI's
clarifying questions, not via AI-written fixes.

\item \textbf{Label encoding.}
I initially planned to encode labels as one-hot vectors
(\texttt{np.zeros(n)} with a single 1.0). The AI asked me to check this against the spec
text ("assigns an integer label") and against how \texttt{nn.CrossEntropyLoss} expects its
\texttt{target} argument (class index, not one-hot). I changed my design to a single
integer label based on this.

\item \textbf{PIL / numpy / tensor conversion.}
I asked for the general syntax to open an image with \texttt{PIL.Image} and convert to a
numpy array, and about channel ordering (HWC vs CHW). The AI explained
\texttt{Image.open}, \texttt{.convert("RGB")}, \texttt{np.array(img)}, and pointed out that
\texttt{torchvision.transforms.ToTensor()} already produces a channel-first tensor,
without writing my \texttt{\_\_getitem\_\_} implementation. I validated this myself by
printing \texttt{img\_tensor.shape} against a shape-mismatch test failure
("314 != 3") I had hit.

\item \textbf{Mean/std computation strategy.}
I asked how to compute per-channel standard deviation at dataset scale. The AI explained
the two equivalent formulas ($\sigma = \sqrt{E[(X-\mu)^2]}$ vs
$\sigma = \sqrt{E[X^2] - (E[X])^2}$) and the single-pass vs two-pass tradeoff, without
giving me the actual mean/std values for this dataset (those I computed myself and will
report separately per the assignment's rounding requirement).

\item \textbf{Train/val split and data leakage.}
I asked whether using the full dataset vs. the train split to compute normalization
statistics mattered. The AI asked me to reconsider computing stats on \texttt{full\_dataset}
before splitting versus on \texttt{train\_subset} after splitting via
\texttt{torch.utils.data.random\_split}. I had initially concluded leakage was
"inevitable" by computing on the full dataset; after this discussion I changed my design
to compute statistics only from the training subset, after splitting.

\item \textbf{\texttt{Dataset}/\texttt{DataLoader} API mechanics.}
General, non-assignment-specific questions about what \texttt{Dataset} provides
(map-style protocol, \texttt{\_\_len\_\_}/\texttt{\_\_getitem\_\_} contract), how
\texttt{DataLoader} batches and collates samples, and how \texttt{random\_split}/
\texttt{Subset} work. These were explained generically (no code specific to my
\texttt{SUN397Dataset} or class folders was written by the AI).

\item \textbf{Validation performed beyond autograder scoring.}
For each item above I verified correctness myself by: re-reading the relevant spec
paragraph and checking my design against its exact wording; running my own code against
the local test failures/tracebacks and inspecting the actual error messages/shapes
(e.g. printing \texttt{kwargs}, printing tensor shapes); and reasoning through small
examples (e.g. what \texttt{os.walk} yields) before writing the corresponding code myself.

\end{itemize}

% Requires \usepackage{enumitem} if you want to keep the [leftmargin=*] styling below;
% otherwise just remove that option.

\begin{itemize}[leftmargin=*]

\item \textbf{Directory traversal / dataset indexing.}
I asked how to enumerate the per-class image folders under \texttt{hw4-data/} and build
an index of (path, label) pairs for \texttt{SUN397Dataset}. The AI pointed me to
\texttt{os.listdir}/\texttt{os.walk}/\texttt{os.path.join}/\texttt{sorted()} conceptually
but did not write the loop. I initially used \texttt{os.walk(dir\_path)} without consuming
it as a generator (tried to unpack it directly into three names), which the AI caught by
asking me to reconsider what a generator yields. I also hit a \texttt{NotADirectoryError}
from \texttt{.DS\_Store} and from checking \texttt{os.path.isfile(dir)} on the wrong
(un-joined) path; both were my bugs, caught via tracebacks I pasted and the AI's
clarifying questions, not via AI-written fixes.

\item \textbf{Label encoding.}
I initially planned to encode labels as one-hot vectors
(\texttt{np.zeros(n)} with a single 1.0). The AI asked me to check this against the spec
text ("assigns an integer label") and against how \texttt{nn.CrossEntropyLoss} expects its
\texttt{target} argument (class index, not one-hot). I changed my design to a single
integer label based on this.

\item \textbf{PIL / numpy / tensor conversion.}
I asked for the general syntax to open an image with \texttt{PIL.Image} and convert to a
numpy array, and about channel ordering (HWC vs CHW). The AI explained
\texttt{Image.open}, \texttt{.convert("RGB")}, \texttt{np.array(img)}, and pointed out that
\texttt{torchvision.transforms.ToTensor()} already produces a channel-first tensor,
without writing my \texttt{\_\_getitem\_\_} implementation. I validated this myself by
printing \texttt{img\_tensor.shape} against a shape-mismatch test failure
("314 != 3") I had hit.

\item \textbf{Mean/std computation strategy.}
I asked how to compute per-channel standard deviation at dataset scale. The AI explained
the two equivalent formulas ($\sigma = \sqrt{E[(X-\mu)^2]}$ vs
$\sigma = \sqrt{E[X^2] - (E[X])^2}$) and the single-pass vs two-pass tradeoff, without
giving me the actual mean/std values for this dataset (those I computed myself and will
report separately per the assignment's rounding requirement).

\item \textbf{Train/val split and data leakage.}
I asked whether using the full dataset vs. the train split to compute normalization
statistics mattered. The AI asked me to reconsider computing stats on \texttt{full\_dataset}
before splitting versus on \texttt{train\_subset} after splitting via
\texttt{torch.utils.data.random\_split}. I had initially concluded leakage was
"inevitable" by computing on the full dataset; after this discussion I changed my design
to compute statistics only from the training subset, after splitting.

\item \textbf{\texttt{Dataset}/\texttt{DataLoader} API mechanics.}
General, non-assignment-specific questions about what \texttt{Dataset} provides
(map-style protocol, \texttt{\_\_len\_\_}/\texttt{\_\_getitem\_\_} contract), how
\texttt{DataLoader} batches and collates samples, and how \texttt{random\_split}/
\texttt{Subset} work. These were explained generically (no code specific to my
\texttt{SUN397Dataset} or class folders was written by the AI).

\item \textbf{Validation performed beyond autograder scoring.}
For each item above I verified correctness myself by: re-reading the relevant spec
paragraph and checking my design against its exact wording; running my own code against
the local test failures/tracebacks and inspecting the actual error messages/shapes
(e.g. printing \texttt{kwargs}, printing tensor shapes); and reasoning through small
examples (e.g. what \texttt{os.walk} yields) before writing the corresponding code myself.

\end{itemize}

% Requires \usepackage{enumitem} if you want to keep the [leftmargin=*] styling below;
% otherwise just remove that option.

\begin{itemize}[leftmargin=*]

\item \textbf{Directory traversal / dataset indexing.}
I asked how to enumerate the per-class image folders under \texttt{hw4-data/} and build
an index of (path, label) pairs for \texttt{SUN397Dataset}. The AI pointed me to
\texttt{os.listdir}/\texttt{os.walk}/\texttt{os.path.join}/\texttt{sorted()} conceptually
but did not write the loop. I initially used \texttt{os.walk(dir\_path)} without consuming
it as a generator (tried to unpack it directly into three names), which the AI caught by
asking me to reconsider what a generator yields. I also hit a \texttt{NotADirectoryError}
from \texttt{.DS\_Store} and from checking \texttt{os.path.isfile(dir)} on the wrong
(un-joined) path; both were my bugs, caught via tracebacks I pasted and the AI's
clarifying questions, not via AI-written fixes.

\item \textbf{Label encoding.}
I initially planned to encode labels as one-hot vectors
(\texttt{np.zeros(n)} with a single 1.0). The AI asked me to check this against the spec
text ("assigns an integer label") and against how \texttt{nn.CrossEntropyLoss} expects its
\texttt{target} argument (class index, not one-hot). I changed my design to a single
integer label based on this.

\item \textbf{PIL / numpy / tensor conversion.}
I asked for the general syntax to open an image with \texttt{PIL.Image} and convert to a
numpy array, and about channel ordering (HWC vs CHW). The AI explained
\texttt{Image.open}, \texttt{.convert("RGB")}, \texttt{np.array(img)}, and pointed out that
\texttt{torchvision.transforms.ToTensor()} already produces a channel-first tensor,
without writing my \texttt{\_\_getitem\_\_} implementation. I validated this myself by
printing \texttt{img\_tensor.shape} against a shape-mismatch test failure
("314 != 3") I had hit.

\item \textbf{Mean/std computation strategy.}
I asked how to compute per-channel standard deviation at dataset scale. The AI explained
the two equivalent formulas ($\sigma = \sqrt{E[(X-\mu)^2]}$ vs
$\sigma = \sqrt{E[X^2] - (E[X])^2}$) and the single-pass vs two-pass tradeoff, without
giving me the actual mean/std values for this dataset (those I computed myself and will
report separately per the assignment's rounding requirement).

\item \textbf{Train/val split and data leakage.}
I asked whether using the full dataset vs. the train split to compute normalization
statistics mattered. The AI asked me to reconsider computing stats on \texttt{full\_dataset}
before splitting versus on \texttt{train\_subset} after splitting via
\texttt{torch.utils.data.random\_split}. I had initially concluded leakage was
"inevitable" by computing on the full dataset; after this discussion I changed my design
to compute statistics only from the training subset, after splitting.

\item \textbf{\texttt{Dataset}/\texttt{DataLoader} API mechanics.}
General, non-assignment-specific questions about what \texttt{Dataset} provides
(map-style protocol, \texttt{\_\_len\_\_}/\texttt{\_\_getitem\_\_} contract), how
\texttt{DataLoader} batches and collates samples, and how \texttt{random\_split}/
\texttt{Subset} work. These were explained generically (no code specific to my
\texttt{SUN397Dataset} or class folders was written by the AI).

\item \textbf{Validation performed beyond autograder scoring.}
For each item above I verified correctness myself by: re-reading the relevant spec
paragraph and checking my design against its exact wording; running my own code against
the local test failures/tracebacks and inspecting the actual error messages/shapes
(e.g. printing \texttt{kwargs}, printing tensor shapes); and reasoning through small
examples (e.g. what \texttt{os.walk} yields) before writing the corresponding code myself.

\end{itemize}

% Requires \usepackage{enumitem} if you want to keep the [leftmargin=*] styling below;
% otherwise just remove that option.

\begin{itemize}[leftmargin=*]

\item \textbf{Directory traversal / dataset indexing.}
I asked how to enumerate the per-class image folders under \texttt{hw4-data/} and build
an index of (path, label) pairs for \texttt{SUN397Dataset}. The AI pointed me to
\texttt{os.listdir}/\texttt{os.walk}/\texttt{os.path.join}/\texttt{sorted()} conceptually
but did not write the loop. I initially used \texttt{os.walk(dir\_path)} without consuming
it as a generator (tried to unpack it directly into three names), which the AI caught by
asking me to reconsider what a generator yields. I also hit a \texttt{NotADirectoryError}
from \texttt{.DS\_Store} and from checking \texttt{os.path.isfile(dir)} on the wrong
(un-joined) path; both were my bugs, caught via tracebacks I pasted and the AI's
clarifying questions, not via AI-written fixes.

\item \textbf{Label encoding.}
I initially planned to encode labels as one-hot vectors
(\texttt{np.zeros(n)} with a single 1.0). The AI asked me to check this against the spec
text ("assigns an integer label") and against how \texttt{nn.CrossEntropyLoss} expects its
\texttt{target} argument (class index, not one-hot). I changed my design to a single
integer label based on this.

\item \textbf{PIL / numpy / tensor conversion.}
I asked for the general syntax to open an image with \texttt{PIL.Image} and convert to a
numpy array, and about channel ordering (HWC vs CHW). The AI explained
\texttt{Image.open}, \texttt{.convert("RGB")}, \texttt{np.array(img)}, and pointed out that
\texttt{torchvision.transforms.ToTensor()} already produces a channel-first tensor,
without writing my \texttt{\_\_getitem\_\_} implementation. I validated this myself by
printing \texttt{img\_tensor.shape} against a shape-mismatch test failure
("314 != 3") I had hit.

\item \textbf{Mean/std computation strategy.}
I asked how to compute per-channel standard deviation at dataset scale. The AI explained
the two equivalent formulas ($\sigma = \sqrt{E[(X-\mu)^2]}$ vs
$\sigma = \sqrt{E[X^2] - (E[X])^2}$) and the single-pass vs two-pass tradeoff, without
giving me the actual mean/std values for this dataset (those I computed myself and will
report separately per the assignment's rounding requirement).

\item \textbf{Train/val split and data leakage.}
I asked whether using the full dataset vs. the train split to compute normalization
statistics mattered. The AI asked me to reconsider computing stats on \texttt{full\_dataset}
before splitting versus on \texttt{train\_subset} after splitting via
\texttt{torch.utils.data.random\_split}. I had initially concluded leakage was
"inevitable" by computing on the full dataset; after this discussion I changed my design
to compute statistics only from the training subset, after splitting.

\item \textbf{\texttt{Dataset}/\texttt{DataLoader} API mechanics.}
General, non-assignment-specific questions about what \texttt{Dataset} provides
(map-style protocol, \texttt{\_\_len\_\_}/\texttt{\_\_getitem\_\_} contract), how
\texttt{DataLoader} batches and collates samples, and how \texttt{random\_split}/
\texttt{Subset} work. These were explained generically (no code specific to my
\texttt{SUN397Dataset} or class folders was written by the AI).

\item \textbf{Validation performed beyond autograder scoring.}
For each item above I verified correctness myself by: re-reading the relevant spec
paragraph and checking my design against its exact wording; running my own code against
the local test failures/tracebacks and inspecting the actual error messages/shapes
(e.g. printing \texttt{kwargs}, printing tensor shapes); and reasoning through small
examples (e.g. what \texttt{os.walk} yields) before writing the corresponding code myself.

\end{itemize}
